%=========================================================================
%  Analyse des données — L3 Économie — CY Cergy Paris Université
%  Séance 1 : Environnement et anatomie des données
%
%  Compilation :  pdflatex 01-intro.tex   (deux fois)
%=========================================================================
\documentclass[10pt,aspectratio=169]{beamer}

\usetheme[progressbar=frametitle,numbering=fraction,block=fill]{metropolis}

\usepackage[T1]{fontenc}
\usepackage[utf8]{inputenc}
\usepackage[french]{babel}
%\usepackage{lmodern}  % non installé ici ; décommentez sur votre machine
\usepackage{booktabs}
\usepackage{listings}
\usepackage{xcolor}

%-------------------------------------------------------------------------
%  PALETTE — remplacer par les valeurs de votre charte CY habituelle
%-------------------------------------------------------------------------
\definecolor{CYprincipal}{HTML}{1B2A4A}   % couleur dominante (titres, barres)
\definecolor{CYaccent}{HTML}{C8102E}      % couleur d'accentuation
\definecolor{CYgris}{HTML}{55637C}        % texte secondaire
\definecolor{CYclair}{HTML}{EEF1F6}       % fonds de blocs

\setbeamercolor{palette primary}{bg=CYprincipal,fg=white}
\setbeamercolor{progress bar}{fg=CYaccent}
\setbeamercolor{frametitle}{bg=CYprincipal,fg=white}
\setbeamercolor{alerted text}{fg=CYaccent}
\setbeamercolor{block title}{bg=CYclair,fg=CYprincipal}
\setbeamercolor{block body}{bg=CYclair!60,fg=black}

%-------------------------------------------------------------------------
%  Mise en forme du code Python
%-------------------------------------------------------------------------
\lstdefinestyle{py}{
  language=Python,
  basicstyle=\ttfamily\small,
  keywordstyle=\color{CYaccent},
  stringstyle=\color{CYprincipal},
  commentstyle=\color{CYgris}\itshape,
  showstringspaces=false,
  breaklines=true,
  frame=none,
  literate={é}{{\'e}}1 {è}{{\`e}}1 {à}{{\`a}}1 {ç}{{\c c}}1 {ê}{{\^e}}1 {ô}{{\^o}}1 {û}{{\^u}}1 {î}{{\^i}}1
}
\lstset{style=py}

%-------------------------------------------------------------------------
\title{Environnement et anatomie des données}
\subtitle{Analyse des données --- Séance 1}
\author{Stefania Marcassa}
\institute{CY Cergy Paris Université --- L3 Économie}
\date{}

\begin{document}

\maketitle

%=========================================================================
\section{Le cours}
%=========================================================================

\begin{frame}{Ce que vous saurez faire en décembre}
\begin{enumerate}
  \item \textbf{Trouver et charger} des données économiques réelles, et lire la documentation d'une source.
  \item \textbf{Nettoyer et restructurer} un jeu de données brut.
  \item \textbf{Produire des statistiques descriptives défendables}.
  \item \textbf{Construire des graphiques honnêtes}.
  \item \textbf{Estimer et lire une régression} en pratique.
  \item \textbf{Rendre un travail reproductible}.
\end{enumerate}
\vfill
\begin{alertblock}{Le critère}
Qu'une autre personne puisse ré-exécuter votre code et retrouver vos résultats.
\end{alertblock}
\end{frame}

\begin{frame}{Ce que ce cours n'est pas}
Ce n'est pas un cours d'économétrie théorique. Les propriétés des estimateurs sont supposées connues.

\vspace{1em}
C'est un cours de \alert{mise en œuvre} : la difficulté est dans les données, pas dans les formules.

\vspace{1.5em}
\begin{block}{Une mise en garde, dès aujourd'hui}
Une régression ne mesure pas un effet causal. Elle mesure une association conditionnelle aux variables incluses.

\vspace{0.4em}
Tout le métier tient dans l'écart entre les deux. Nous y reviendrons à chaque séance à partir de la séance~7.
\end{block}
\end{frame}

\begin{frame}{Comment le semestre est organisé}
\begin{columns}[T]
\begin{column}{0.48\textwidth}
\textbf{Trois blocs}
\begin{itemize}
  \item Séances 1--5 : manipulation des données
  \item Séances 6--8 : visualisation et régression
  \item Séances 9--10 : séries temporelles, synthèse
\end{itemize}
\end{column}
\begin{column}{0.48\textwidth}
\textbf{Chaque TD}
\begin{itemize}
  \item 25 min : retour sur l'exercice, erreurs fréquentes
  \item 30 min : résolution en direct
  \item 45 min : travail autonome
\end{itemize}
\end{column}
\end{columns}

\vspace{1.2em}
Chaque séance donne lieu à trois fichiers sur le site du cours : le \textbf{notebook à trous} (avant), le \textbf{script complet} (après), l'\textbf{exercice} (pour la fois suivante).
\end{frame}

\begin{frame}{Évaluation}
\begin{center}
\begin{tabular}{lccl}
\toprule
\textbf{Épreuve} & \textbf{Poids} & \textbf{Durée} & \textbf{Périmètre} \\
\midrule
Contrôle continu & 50\,\% & 1 h & Séances 1 à 5 \\
Contrôle terminal & 50\,\% & 2 h & Ensemble du cours \\
\bottomrule
\end{tabular}
\end{center}

\vspace{1.2em}
Les deux épreuves se déroulent \alert{sur papier, sans machine}.

\vspace{0.8em}
Aucune ne demande d'écrire du code de mémoire.
\end{frame}

\begin{frame}{Quatre types de questions}
\begin{enumerate}
  \item \textbf{Lecture de sortie.} Un tableau ou une régression vous est donné : que montre-t-il, et que ne montre-t-il pas ?
  \item \textbf{Diagnostic.} Un extrait de données ou de code contient un problème : l'identifier, proposer une correction.
  \item \textbf{Interprétation économique.} Que peut-on affirmer à partir de ce coefficient, et sous quelles conditions ?
  \item \textbf{Code court.} Trois à cinq lignes à écrire ou compléter --- jamais un script entier.
\end{enumerate}
\vfill
Ce qui est évalué n'est pas la syntaxe. C'est le \alert{jugement sur les données}.
\end{frame}

\begin{frame}{Les assistants d'IA}
Autorisés et attendus pour les exercices. C'est l'outil de travail réel de votre future profession.

Interdits lors des deux épreuves, qui se tiennent sans machine.

\vspace{1.2em}
\begin{block}{Ce qu'ils font bien, ce qu'ils font mal}
Un assistant écrit très bien la syntaxe.

\vspace{0.4em}
Il est peu fiable sur les décisions qui font l'objet de ce cours : quelle variable retenir, comment traiter les manquants, quelle spécification est défendable, ce qu'un coefficient autorise à conclure.
\end{block}

\vspace{0.8em}
Si vous ne pouvez pas expliquer une ligne qu'il a produite, vous ne saurez pas répondre à la question de diagnostic qui portera dessus.
\end{frame}

%=========================================================================
\section{L'environnement de travail}
%=========================================================================

\begin{frame}{Google Colab, et rien d'autre}
Tout le cours se fait dans le navigateur.

\vspace{1em}
\begin{itemize}
  \item Aucune installation, aucune version de Python à gérer
  \item Aucun problème de chemin d'accès
  \item Un compte Google suffit
  \item Le code s'exécute sur une machine distante, pas sur la vôtre
\end{itemize}

\vspace{1.2em}
Le travail en local (Anaconda, VS Code) est possible mais \alert{non assisté} : les TD ne sont pas consacrés au dépannage d'installations.
\end{frame}

\begin{frame}{Premier notebook --- à faire maintenant}
\begin{enumerate}
  \item Connectez-vous à votre compte Google.
  \item Allez sur \texttt{colab.research.google.com}
  \item Cliquez sur \textbf{Nouveau notebook}.
  \item Utilisez Chrome ou Firefox de préférence.
\end{enumerate}

\vspace{1.2em}
Un notebook est une suite de \textbf{cellules} de deux types :
\begin{itemize}
  \item les cellules de \textbf{code}, que l'on exécute avec \texttt{Maj + Entrée} ;
  \item les cellules de \textbf{texte}, pour commenter ce que l'on fait.
\end{itemize}

\vspace{0.8em}
Un notebook sans cellules de texte est un notebook illisible --- y compris pour vous, dans trois semaines.
\end{frame}

\begin{frame}[fragile]{Une première cellule}
\begin{lstlisting}
# Le taux de croissance entre deux annees
pib_2023 = 2822
pib_2024 = 2891

croissance = (pib_2024 - pib_2023) / pib_2023
print(croissance)
\end{lstlisting}

\vspace{0.5em}
\begin{lstlisting}
0.024450035435861798
\end{lstlisting}

\vspace{1em}
Trois choses à retenir : le \texttt{\#} introduit un commentaire, le \texttt{=} affecte une valeur, et rien ne s'affiche sans \texttt{print()}.
\end{frame}

\begin{frame}{Un piège qui vous coûtera dix minutes}
Lors du premier import de fichier, Colab peut renvoyer une erreur si les \textbf{cookies tiers} sont bloqués dans votre navigateur.

\vspace{1.2em}
La correction :
\begin{enumerate}
  \item Ouvrir les paramètres du navigateur
  \item Autoriser \texttt{https://[*.]googleusercontent.com:443}
\end{enumerate}

\vspace{1.2em}
Notez-le maintenant. Vous en aurez besoin à la séance 2.
\end{frame}

%=========================================================================
\section{Les objets dont nous aurons besoin}
%=========================================================================

\begin{frame}[fragile]{Quatre types de base}
\begin{lstlisting}
inflation = 2.4          # float  : nombre a virgule
annee     = 2024         # int    : nombre entier
pays      = "France"     # str    : texte
zone_euro = True         # bool   : vrai / faux
\end{lstlisting}

\vspace{1.2em}
La distinction \texttt{int} / \texttt{float} paraît anodine. Elle ne l'est pas :
\begin{itemize}
  \item un identifiant de commune lu comme un \texttt{int} perd son zéro initial ;
  \item une année lue comme un \texttt{str} ne se compare pas, ne se trie pas.
\end{itemize}

\vspace{0.6em}
La moitié des bugs de ce cours viendra d'un type mal interprété à la lecture d'un fichier.
\end{frame}

\begin{frame}[fragile]{Listes : une suite ordonnée}
\begin{lstlisting}
prix = [1.20, 0.95, 2.40, 1.75]

prix[0]        # 1.2   -- on compte a partir de zero
prix[-1]       # 1.75  -- le dernier
len(prix)      # 4
sum(prix)      # 6.3
\end{lstlisting}

\vspace{1.2em}
L'indexation commence à \alert{zéro}. C'est la source d'erreur la plus fréquente chez les débutants, et elle ne produit pas de message d'erreur : elle produit un résultat faux.
\end{frame}

\begin{frame}[fragile]{Dictionnaires : des paires clé--valeur}
\begin{lstlisting}
panier = {"pain": 1.20, "lait": 0.95, "cafe": 2.40}

panier["pain"]          # 1.2
panier.keys()           # les biens
panier.values()         # les prix
\end{lstlisting}

\vspace{1.2em}
Un dictionnaire décrit exactement ce dont nous avons besoin en TD : un \textbf{panier de biens}, où chaque bien porte un prix et une quantité.

\vspace{0.8em}
C'est aussi, en pratique, la structure d'une ligne d'un tableau de données.
\end{frame}

%=========================================================================
\section{Anatomie d'un tableau de données}
%=========================================================================

\begin{frame}{Trois mots à ne jamais confondre}
\begin{center}
\begin{tabular}{lll}
\toprule
\textbf{id\_menage} & \textbf{revenu} & \textbf{region} \\
\midrule
1001 & 32\,400 & Île-de-France \\
1002 & 28\,900 & Occitanie \\
1003 & 41\,200 & Bretagne \\
\bottomrule
\end{tabular}
\end{center}

\vspace{1.5em}
\begin{itemize}
  \item Une \textbf{observation} : une ligne. Ici, un ménage.
  \item Une \textbf{variable} : une colonne. Ici, le revenu.
  \item L'\textbf{unité statistique} : ce que représente une ligne. Ici, le ménage --- et non l'individu.
\end{itemize}
\end{frame}

\begin{frame}{Pourquoi l'unité statistique est le premier piège}
La même question économique change de réponse selon l'unité retenue.

\vspace{1.2em}
\begin{block}{Un exemple}
« Quel est le revenu moyen ? »

\vspace{0.5em}
Par \textbf{ménage} : un chiffre. Par \textbf{individu} : un autre. Par \textbf{unité de consommation} : un troisième.

\vspace{0.5em}
Aucun n'est faux. Ils ne répondent simplement pas à la même question.
\end{block}

\vspace{1.2em}
Avant toute analyse : \alert{de quoi une ligne est-elle la description ?}
\end{frame}

\begin{frame}{Trois structures de données}
\begin{columns}[T]
\begin{column}{0.31\textwidth}
\textbf{Coupe transversale}

\vspace{0.5em}
Plusieurs unités, une date.

\vspace{0.5em}
\small\color{CYgris}
Les salaires de 40\,000 individus en 2023.
\end{column}
\begin{column}{0.31\textwidth}
\textbf{Série temporelle}

\vspace{0.5em}
Une unité, plusieurs dates.

\vspace{0.5em}
\small\color{CYgris}
Le taux de chômage français, 1975--2025.
\end{column}
\begin{column}{0.31\textwidth}
\textbf{Panel}

\vspace{0.5em}
Plusieurs unités, plusieurs dates.

\vspace{0.5em}
\small\color{CYgris}
Le PIB de 27 pays, 2000--2024.
\end{column}
\end{columns}

\vspace{2em}
La structure détermine ce que l'on peut estimer --- et ce que l'on ne peut pas. Nous construirons un panel à la séance~4.
\end{frame}

\begin{frame}{Exercice éclair}
Pour chacun de ces jeux de données : quelle est l'unité statistique, et quelle est la structure ?

\vspace{1.5em}
\begin{enumerate}
  \item Le prix de clôture du CAC 40, chaque jour depuis 1990.
  \item Les réponses à l'Enquête Emploi du deuxième trimestre 2024.
  \item Les transactions immobilières enregistrées en France de 2014 à 2024.
  \item Le taux d'emploi des 27 pays de l'UE, chaque année depuis 2000.
\end{enumerate}

\vspace{1.5em}
\small\color{CYgris}
Le cas 3 est moins évident qu'il n'y paraît : discutons-en.
\end{frame}

%=========================================================================
\section{Pour la suite}
%=========================================================================

\begin{frame}{TD 1 --- Un indice des prix}
Nous construirons un indice des prix à la consommation à partir d'un panier de biens pondéré, puis nous comparerons deux façons de le calculer.

\vspace{1.2em}
\begin{itemize}
  \item L'indice de \textbf{Laspeyres} fixe les quantités à l'année de base.
  \item L'indice de \textbf{Paasche} les fixe à l'année courante.
\end{itemize}

\vspace{1.2em}
Les deux mesurent « l'inflation ». Ils ne donnent pas le même chiffre.

\vspace{0.8em}
Toute la question du TD : \alert{pourquoi}, et lequel choisir.
\end{frame}

\begin{frame}{Avant la séance 2}
\begin{enumerate}
  \item Créez votre compte Google si vous n'en avez pas.
  \item Ouvrez un notebook Colab et exécutez une cellule.
  \item Autorisez \texttt{googleusercontent.com} dans votre navigateur.
  \item Récupérez le notebook de la séance 2 sur le site du cours.
\end{enumerate}

\vspace{2em}
\begin{center}
\large \texttt{stefaniamarcassa.github.io/analyse\_des\_donnees}
\end{center}
\end{frame}

\begin{frame}[standout]
Des questions ?
\end{frame}

\end{document}